\documentclass[11pt]{article}

\usepackage[margin=1in]{geometry}
\usepackage{amsmath}
\usepackage{amssymb}
\usepackage{amsthm}
\usepackage{booktabs}
\usepackage{hyperref}

\newtheorem{definition}{Definition}
\newtheorem{theorem}{Theorem}
\newtheorem{lemma}{Lemma}
\newtheorem{corollary}{Corollary}
\theoremstyle{remark}
\newtheorem{remark}{Remark}

\title{Conjecture 00000005626 is False:\\ a Six-Point Set in General Position with No Empty Pentagon}
\author{earthking11}
\date{2026-09-13}

\begin{document}
\maketitle

\begin{abstract}
We refute conjecture 00000005626, which asserts that every $6$-point subset
of the plane in general position contains an empty pentagon.  We exhibit an
explicit six-point set in general position whose number of empty pentagons
is zero.  The example is a convex pentagon with one strictly interior point,
with all coordinates integral so that every verification is exact.  We also
note that the conjecture's additional clause about a ``unique Horton-type
counterexample family'' is incoherent, and recall that the correct threshold
for the empty-pentagon problem is $10$ points ($9$-point Horton sets are
extremal).
\end{abstract}

\section{The conjecture}

Conjecture 00000005626 (translated) concerns the empty-polygon spectrum of
Horton-type recursive constructions and includes the following two claims:

\begin{enumerate}
\item[(C1)] every $6$-point subset of the plane in general position contains
an empty pentagon;
\item[(C2)] Horton type is the unique counterexample family.
\end{enumerate}

We show that (C1) is false by producing a six-point set in general position
that contains no empty pentagon.

\section{Empty pentagons}

\begin{definition}[empty pentagon]
A $5$-element subset $S$ of a finite point set $X$ is an \emph{empty
pentagon} if
\begin{enumerate}
\item[(i)] $S$ is in \emph{strictly convex position}, i.e.\ every one of its
five points is an extreme point of $\operatorname{conv}(S)$; and
\item[(ii)] $\operatorname{conv}(S)$ contains no other point of $X$.
\end{enumerate}
\end{definition}

Thus a $5$-subset fails to be an empty pentagon if some of its points is
interior to the hull of the other four (or lies on an edge, so is not an
extreme point), or if the one omitted point of $X$ lies in
$\operatorname{conv}(S)$.

\section{The counterexample}

Let
\begin{equation}
P_1=(0,0),\quad P_2=(4,0),\quad P_3=(5,2),\quad P_4=(2,4),\quad P_5=(-1,2),
\end{equation}
and let
\begin{equation}
C=(2,1).
\end{equation}
Let $X=\{P_1,\dots,P_5,C\}$.

For points $o,a,b\in\mathbb{Z}^2$ define the oriented cross product
\begin{equation}
\operatorname{cross}(o,a,b)
  =(a_1-o_1)(b_2-o_2)-(a_2-o_2)(b_1-o_1)\in\mathbb{Z}.
\end{equation}
It is positive when $o\to a\to b$ is a left turn, negative for a right turn,
and zero when the three points are collinear.

\begin{remark}[why not the obvious centre $(2,2)$]
The lattice point $(2,2)$ is collinear with $P_3=(5,2)$ and $P_5=(-1,2)$,
since all three have second coordinate $2$.  Hence $\{P_1,\dots,P_5,(2,2)\}$
is \emph{not} in general position, and cannot serve as a counterexample to
(C1).  Among all lattice points strictly interior to the pentagon,
$C=(2,1)$ is the unique one that is strictly inside the hull of every four
of the five outer vertices and keeps the six-point set in general position;
this was verified by exhaustive finite search (see \texttt{reproduce.py}).
\end{remark}

\section{Verification}

\subsection{General position}

\begin{lemma}
No three points of $X$ are collinear, i.e.\ $X$ is in general position.
\end{lemma}

\begin{proof}
There are $\binom{6}{3}=20$ triples.  Direct exact integer computation gives
$\operatorname{cross}\neq 0$ for all of them.  Equivalently, the Lean theorem
\texttt{general\_position} proves this by kernel evaluation.
\end{proof}

\subsection{The outer five points are strictly convex}

The five consecutive cross products of the outer pentagon are
\begin{equation}
\begin{aligned}
\operatorname{cross}(P_1,P_2,P_3)&=8, &
\operatorname{cross}(P_2,P_3,P_4)&=8, &
\operatorname{cross}(P_3,P_4,P_5)&=12,\\
\operatorname{cross}(P_4,P_5,P_1)&=8, &
\operatorname{cross}(P_5,P_1,P_2)&=8. &&
\end{aligned}
\end{equation}
All are strictly positive, so $P_1P_2P_3P_4P_5$ is a strictly convex
counter-clockwise pentagon.

\subsection{The centre is strictly interior}

For a counter-clockwise convex polygon, a point lies strictly inside iff it
is strictly to the left of every directed edge.  For $C=(2,1)$ we obtain
\begin{equation}
\operatorname{cross}(P_1,P_2,C)=4,\quad
\operatorname{cross}(P_2,P_3,C)=5,\quad
\operatorname{cross}(P_3,P_4,C)=9,
\end{equation}
\begin{equation}
\operatorname{cross}(P_4,P_5,C)=9,\quad
\operatorname{cross}(P_5,P_1,C)=5.
\end{equation}
All are strictly positive, so $C$ is strictly inside
$\operatorname{conv}(P_1,\dots,P_5)$.  In fact, checking the four edges of
each of the five quadrilaterals obtained by deleting one outer vertex shows
that $C$ is strictly inside the hull of \emph{every} four of the five outer
vertices.

\subsection{There is no empty pentagon}

The six $5$-subsets of $X$ are as follows.

\begin{center}
\begin{tabular}{lll}
\toprule
$5$-subset & omitted & why it is not an empty pentagon\\
\midrule
$\{P_1,P_2,P_3,P_4,P_5\}$ & $C$ & strictly convex, but $C\in\operatorname{conv}(\cdot)$\\
$\{P_2,P_3,P_4,P_5,C\}$ & $P_1$ & $C$ non-extreme, not convex position\\
$\{P_1,P_3,P_4,P_5,C\}$ & $P_2$ & $C$ non-extreme, not convex position\\
$\{P_1,P_2,P_4,P_5,C\}$ & $P_3$ & $C$ non-extreme, not convex position\\
$\{P_1,P_2,P_3,P_5,C\}$ & $P_4$ & $C$ non-extreme, not convex position\\
$\{P_1,P_2,P_3,P_4,C\}$ & $P_5$ & $C$ non-extreme, not convex position\\
\bottomrule
\end{tabular}
\end{center}

Indeed, in the first subset the hull of the five outer points contains the
interior point $C$, so condition (ii) fails.  In each of the remaining five
subsets the point $C$ lies strictly inside the hull of the four outer
vertices present, hence is not an extreme point; condition (i) fails.

\begin{theorem}
The set $X$ is in general position and contains zero empty pentagons.
Consequently claim (C1) of conjecture 00000005626 is false.
\end{theorem}

\begin{proof}
Immediate from the three lemmas above: enumerate the six $5$-subsets as in
the table.  For each, either strict convexity or emptiness fails.  Hence the
number of empty pentagons is $0$.
\end{proof}

This generalises: an $(n-1)$-point strictly convex polygon together with a
single strictly interior point has no empty $n$-gon, because the interior
point lies in the hull of every $n$-subset omitting it, while in any
$n$-subset containing it, it is non-extreme.

\section{The second clause is incoherent}

Clause (C2) says that Horton type is the unique counterexample family.  But
(C1) asserts that \emph{every} general-position $6$-point set contains an
empty pentagon.  If (C1) were true there would be no counterexample family
at all, so no ``unique counterexample family'' could exist.  The two clauses
cannot both be meaningful.  Our counterexample is not of Horton type: it is
a convex pentagon with one interior point.

\section{The correct threshold is $10$}

The empty-pentagon problem has sharp threshold $10$: there exist
$9$-point sets in general position with no empty pentagon, the classical
extremal examples being nine-point \emph{Horton sets}, while every
$10$-point set in general position in the plane contains an empty pentagon.
Thus the threshold asserted implicitly by (C1) (namely $6$) is far too low;
our six-point example is a small explicit witness.

\section{Formal and computational verification}

All claims above were checked mechanically.

\begin{itemize}
\item \texttt{lean4/} contains a Lean~4 formalisation on toolchain
\texttt{leanprover/lean4:v4.33.1}, using only core Lean and \texttt{Std}
(no Mathlib, no $\mathbb{R}$, no \texttt{sorry}, no \texttt{axiom}, no
\texttt{native\_decide}).  Points are integer pairs; all geometry is exact
\texttt{Int} arithmetic; all proofs are by kernel \texttt{decide}.  The
theorems \texttt{general\_position}, \texttt{pentagon\_strictly\_convex},
\texttt{centre\_strictly\_inside}, \texttt{no\_empty\_pentagon}, and
\texttt{conjecture\_00000005626\_false} each depend on \emph{no axioms}.
\item \texttt{reproduce.py} (standard library only) verifies general
position, strict convexity, strict interiority, and that the empty-pentagon
count over all $\binom{6}{5}=6$ subsets is $0$, using exact integer
arithmetic; it also repeats the count for a regular pentagon plus its centre
using \texttt{sympy} when available.  It prints \texttt{PASS} and exits $0$.
\end{itemize}

\end{document}
